\begin{textsample}{POS Dim 4 – human – Score 4.00 – mg/mg\_ef\_linguagens\_ingles\_9\_p2.txt}  \label{ex:f4_pos_019}
O sentido não é uma característica do \textbf{texto}, mas se constrói no diálogo leitor-texto, numa espécie de jogo psicolinguístico que se estabelece durante o processamento da informação.
Central é a noção de que o \textbf{texto} não é um produto acabado, mas é (re)criado a cada nova \textbf{leitura}.
Skimming : olhada rápida ( um passar de olhos ) pelo \textbf{texto} para ter uma ideia geral do assunto tratado.
Scanning : localização rápida de informação no \textbf{texto}.
Identificação do padrão geral de organização ( gênero \textbf{textual} ).
Uso de pistas não-verbais ( ilustrações, diagramas, tabelas, saliências gráficas, etc. ).
Uso de títulos, subtítulos, legendas, suporte ( ou portador ) do \textbf{texto}.
Antecipações do que vem em seguida ao que está sendo lido.
Uso do contexto e cognatos.
Uso de pistas \textbf{textuais} ( pronomes, conectivos, articuladores, etc. ).
Construção dos elos coesivos ( lexicais e gramaticais ).
Identificação do tipo do \textbf{texto} e das articulações na superfície \textbf{textual}.
Uso de palavras-chave para construir a progressão temática.
Construção de inferências.
Transferência de informação : do verbal para o não-verbal ( resumos do que foi lido na forma de tabelas, esquemas ou mapas conceituais ). ( NUNAN, 1999. p.265-266 ).

% matched lemmas: leitura, texto, textual
\end{textsample}
